\documentclass[10pt,a4paper]{amsart}
\usepackage[margin=19mm]{geometry}
\usepackage{amsmath,amssymb,mathtools}
\usepackage[scr=rsfs,cal=euler]{mathalfa}
\usepackage{xfrac}
\usepackage[unicode,hidelinks,bookmarks=false]{hyperref}
\newcommand{\ie}{i.e.\ }
\newcommand{\cf}{cf.\ }
\newcommand{\resp}{respectively\ }
\newcommand{\Subsection}{\subsection}
\providecommand{\nobreakdash}{\nobreak\hbox{-}}
\setlength{\emergencystretch}{2em}
\multlinegap0pt
\usepackage{bm}
\usepackage{bbm}
\usepackage{dsfont}
\usepackage{xspace}
\usepackage{cancel}
\usepackage{mathtools}
\usepackage{amsthm}
\makeatletter
\@ifundefined{thm}{\newtheorem{thm}{Theorem}}{}
\@ifundefined{defi}{\newtheorem{defi}[thm]{Definition}}{}
\@ifundefined{lemma}{\newtheorem{lemma}[thm]{Lemma}}{}
\@ifundefined{prop}{\newtheorem{prop}[thm]{Proposition}}{}
\makeatother
\newcommand*\dif{\mathop{}\!\mathrm{d}}
\newcommand{\bs}[1]{\boldsymbol{#1}}
\newcommand{\mr}[1]{\mathrm{#1}}
\newcommand{\msc}[1]{\mathscr{#1}}
\newcommand{\wt}[1]{\widetilde{#1}}
\title[On the structure of entropy dissipation and regularity for q]{On the structure of entropy dissipation and regularity for quasi-entropy solutions to 1d scalar conservation laws and to isentropic Euler system with $\gamma=3$}
\author{Fabio Ancona, Elio Marconi, Luca Talamini}
\date{Source: arXiv:2601.03739v1 (CC BY 4.0)}
\begin{document}
\maketitle

\noindent\emph{Source and licence.} On the structure of entropy dissipation and regularity for quasi-entropy solutions to 1d scalar conservation laws and to isentropic Euler system with $\gamma=3$. Version arXiv:2601.03739v1 (CC BY 4.0), distributed under the Creative Commons Attribution 4.0 International licence, as recorded in the arXiv licence metadata for this exact version and shown on the abstract page https://arxiv.org/abs/2601.03739v1. Full rights notice: licenses/arxiv-2601.03739-CC-BY-4.0.txt.\par\smallskip

\noindent\textit{Editorial scope.} Counted unit: the Thm whose source label is \texttt{thm:besov} in the article (source0.tex lines 1461--1467, source label \texttt{thm:besov}), delivered together with the article's own complete original proof of it (source0.tex lines 1470--1602, one \texttt{proof} environment of 133 lines). No mathematics is written, repaired or re-derived: statement and proof are the article's own text line for line. Required context retained uncounted: eq:conslaw (lines 84--86); defi:fes (lines 154--159); prop:feskin (lines 176--189); eq:chidef (lines 183--188); def:lagrangianrep (lines 296--321); eq:charaeqla (lines 308--312); eq:mu1def (lines 360--366); defi:goodlr (lines 396--403); defi:lagraepi (lines 684--696); defi:siminduc (lines 722--724); prop:goodpairs (lines 742--744); lemma:goodcurves (lines 768--780); lemma:nocrossing (lines 1027--1037). Every definition, display and label of those lines is retained unchanged. Omitted from this package and not redistributed: 1--83, 87--153, 160--175, 189--295, 313--359, 367--395, 404--683, 697--721, 725--741, 745--767, 781--1026, 1038--1460, 1468--1469, 1603--2175. Nothing inside a retained excerpt is omitted or altered: every retained body is delivered exactly as the article's lines are written, so no omission receipt of any kind accompanies this package. Citations: the retained excerpts carry no citation command at all, so no bibliography entry of the article is transcribed and no reference is redistributed. Reference adaptation: none was needed and none was made; every reference inside the delivered unit resolves inside it, so no unresolved reference or citation is produced. Printed slips kept literal: no printed word, symbol, hypothesis or number of the counted statement or of its proof was altered. Layout and class: the article's own class is \texttt{amsart} with packages including geometry, xcolor, hyperref, amsmath, enumerate, tikz, comment, bbm, amssymb, mathrsfs, stmaryrd, enumitem, dirtytalk, tikz-cd; the wrapper is the lane's pinned \texttt{amsart} editorial wrapper at 10pt on A4, which restates the theorem-like environments the retained text needs and adds the article's own macro definitions that the retained text needs (\texttt{\textbackslash bs}, \texttt{\textbackslash dif}, \texttt{\textbackslash mr}, \texttt{\textbackslash msc}, \texttt{\textbackslash wt}), with extra wrapper packages (bm, bbm, dsfont, xspace, cancel, mathtools). The delivered numbering is the unit's own. Assets: no retained span contains a figure environment, a graphics-inclusion command, a PDF-page inclusion, a table or a TikZ picture; no image file is copied into this package, the package declares no asset dependency and no image file is redistributed with it. Licence: version arXiv:2601.03739v1 of this article is distributed under the Creative Commons Attribution 4.0 International licence, as recorded in the arXiv licence metadata for this exact version and shown on the abstract page https://arxiv.org/abs/2601.03739v1. A full rights notice is delivered at licenses/arxiv-2601.03739-CC-BY-4.0.txt. Characters: every retained excerpt is pure ASCII, so the delivered engine, XeLaTeX with its default Latin Modern text font, reports no missing character.
\begin{equation}\label{eq:conslaw}
    \partial_ t u + \mr{div}_x f(u) = 0,  \qquad \text{in} \quad (0, T) \times \mathbb R^d
\end{equation}

\begin{defi}\label{defi:fes}
    We say that $u \in \mathbf L^\infty(\Omega; \mathbb R)$ is a \textit{quasi-entropy} solution with flux $f$ if for every entropy-entropy flux pair $(\eta, \bs q)$, with $\eta \in C^2(\mathbb R)$, there exists a locally finite measure $\mu_{\eta} \in \msc M(\Omega)$ such that 
    \begin{equation}\label{eq:fenergy}
        \partial_t \eta(u) + \mr{div}_x \, \bs q(u) \, = \, \mu_{\eta} \qquad \text{in $\msc D^\prime_{t,x}(\Omega)$}.
    \end{equation}
\end{defi}

\begin{prop}\label{prop:feskin}
    A function $u \in \mathbf L^\infty(\Omega; \mathbb R)$ is a quasi-entropy solution in the sense of Definition \ref{defi:fes} if and only if, 
there exist locally finite measures $\mu_0, \mu_1 \in \msc M(\Omega \times \mathbb R)$ with $\mr{supp}\, \mu_i \subset \Omega \times [\inf u, \sup u]$, such that 
\begin{equation}\label{eq:kinu}
    \partial_t \chi + f^\prime(v) \cdot \nabla_x \chi = \partial_v \mu_1 + \mu_0 \qquad \text{in $\msc D^\prime(\Omega \times \mathbb R)$}
\end{equation}
where 
\begin{equation}\label{eq:chidef}
\chi(t, x, v) = \begin{cases}
    1 & \text{if $v \leq u(t,x)$},\\
    0 & \text{otherwise}.
\end{cases}
\end{equation}
\end{prop}

\begin{defi}\label{def:lagrangianrep}
A \textit{Lagrangian representation} (of the hypograph) of a quasi-entropy solution $u$ to \eqref{eq:conslaw} is a locally finite measure $\bs \omega \in \msc {M}^+(\Gamma)$ such that:
\begin{enumerate}
    \item it holds
    \begin{equation}\label{eq:deflagromega}
        \chi \cdot \mathscr L^{d+2} = \int_{\Gamma} (\mathrm{id}, \gamma)_{\sharp} \mathscr L^1 \llcorner I_{\gamma} \dif \bs \omega(\gamma)
    \end{equation}
    where $\chi$ is defined as in \eqref{eq:chidef}, and the measure in the right hand side is defined on test functions $\varphi \in  C_c(\mathbb R^{d+2})$ by
    $$
    \int_{\mathbb R^{d+2}} \varphi \dif \Big[ \int_{\Gamma} (\mathrm{id}, \gamma)_{\sharp} \mathscr L^1 \llcorner I_{\gamma} \dif \bs \omega(\gamma)\Big] := \int_\Gamma \Big( \int_{I_\gamma} \varphi(t,\gamma(t)) \dif t\Big) \dif \bs \omega(\gamma).
    $$
    \item $\omega$ is concentrated on the set of curves $\gamma = (\gamma^x, \gamma^v) \in \Gamma$ such that 
    \begin{equation}\label{eq:charaeqla}
    \begin{aligned}
   &  \dot{\gamma}^x = f^{\prime}(\gamma^v(t)), \qquad \text{for $\msc L^1$-a.e. $t \in I_{\gamma}$}, 
   \end{aligned}
    \end{equation}    
    \item the following integral bound holds: for every compact $K \subset \mathbb R^d$ and $T > 0$, one has
    \begin{equation}\label{eq:CKTlr}
        \begin{aligned}
     \int_{\Gamma}  \mathrm{Tot.Var.}\big(\gamma^v, \, I_\gamma \cap [0,T]\big) \cdot \mathbf 1_{\{\gamma^x(t_\gamma^1) \in K\}}(\gamma) \,\dif \bs \omega(\gamma) < C_{K,T}\\
\sum_{i=1}^2 \bs \omega \Big(\big\{ \gamma \in \Gamma \; | \; (t,\gamma^x(t^i_\gamma)) \in (0, T) \times K \big\} \Big) \leq \wt C_{K,T}
     \end{aligned}
    \end{equation}
\end{enumerate}
\end{defi}

\begin{equation}\label{eq:mu1def}
    \begin{aligned}
    \mu_1^{\gamma}& \doteq -(\mathbb I , \gamma)_{\sharp} \Tilde{D} \gamma^v \\
    & - \mathscr H^1 \llcorner \{(t,x,v) \; | \; \gamma^v(t-) \leq v \leq \gamma^v(t+), \quad t \in J_\gamma\}\\
    & +\mathscr H^1 \llcorner \{(t,x,v) \; | \; \gamma^v(t+) \leq v \leq \gamma^v(t-), \quad t \in J_\gamma\},
    \end{aligned}
 \end{equation}

\begin{defi}\label{defi:goodlr}
We say that a Lagrangian representation $\bs \omega$ is \emph{good} if the pair $(\mu_0, \mu_1) \in \mathscr M(\Omega \times \mathbb R)^2$ induced by $\bs \omega$ satisfies additionally
\begin{equation}\label{eq:optcond}
    \begin{aligned}
        \left| \mu_1 \right| = \int_{\Gamma} |\mu_1^{\gamma}| \dif \bs \omega(\gamma), \qquad \left| \mu_0 \right| = \int_{\Gamma} |\mu_0^{\gamma}| \dif \bs \omega(\gamma).
    \end{aligned}
\end{equation}
\end{defi}

\begin{defi}\label{defi:lagraepi}
We say that $\bs \omega_e$ is a Lagrangian representation for the epigraph of $u$ if conditions (2), (3) of Definition \ref{def:lagrangianrep} hold, and condition (1) is replaced by:
\begin{equation}\label{eq:deflagromegae}
        \chi_e \cdot \mathscr L^{d+2} = \int_{\Gamma} (\mathrm{id}, \gamma)_{\sharp} \mathscr L^1 \llcorner I_{\gamma} \dif \bs \omega_e(\gamma)
    \end{equation}
where 
\begin{equation}
    \chi_e(t,x,v) \doteq \begin{cases}
        1 & \text{if} \; u(t, x) \leq v , \\
        0 & \text{otherwise}
    \end{cases}
\end{equation}
\end{defi}

\begin{defi}\label{defi:siminduc}
    We say that a pair $(\mu_0, \mu_1)$ is \textit{simultaneously induced} by Lagrangian representations $\bs \omega_h$, $\bs \omega_e$ of the hypograph and of the epigraph respectively if $(\mu_0, \mu_1)$ is induced by $\bs \omega_h$ and $(-\mu_0, -\mu_1)$ is induced by $\bs \omega_e$.
\end{defi}

\begin{prop}\label{prop:goodpairs}
    Let $u$ be a quasi-entropy solution. Every pair $(\widehat \mu_0, \widehat \mu_1) \in \mathbf P_\chi$  which is minimal with respect to $\preceq$ is simultaneously induced by some good Lagrangian representations $\bs \omega_h, \bs \omega_e$.
\end{prop}

   \begin{lemma}\label{lemma:goodcurves}
       For $\bs \omega_h$-a.e. $\gamma \in \Gamma$ it holds that for $\msc L^1$-a.e. $t \in [0,T]$
       \begin{enumerate}
           \item $(t, \gamma^x(t))$ is a Lebesgue point of $u$;
           \item $\gamma^v(t) < u(t, \gamma^x(t))$
       \end{enumerate}
       We denote by $\Gamma_h$ the set of curves $\gamma \in \Gamma$ such that the two properties above hold. Similarly, for $\bs \omega_e$-a.e. $\gamma \in \Gamma$ it holds that for $\msc L^1$-a.e. $t \in [0,T]$
       \begin{enumerate}
           \item $(t, \gamma^x(t))$ is a Lebesgue point of $u$;
           \item $\gamma^v(t) > u(t, \gamma^x(t))$
       \end{enumerate}
       and we denote by $\Gamma_e$ the set of curves $\gamma \in \Gamma$ such that the two properties above hold. 
\end{lemma}

\begin{lemma}\label{lemma:nocrossing}
    Let $\Gamma_h, \Gamma_e$ as in Lemma \ref{lemma:goodcurves}. Let $(\bar \gamma, I_{\bar \gamma}) \in \Gamma_h$ and let $(a, b) \subset I_{\bar \gamma}$ 
 be such that $\bar \gamma^v((a, b) )\subset A_l$. Let moreover $G \subset \Gamma_e$ be a set of curves $(\gamma, I_{\gamma})$ such that there exist $a < s^1_{\gamma} < s^2_{\gamma}< b$ with
\begin{equation}\label{eq:crossingcond}
    \gamma^x(s^1_{\gamma}) > \bar \gamma^x (s^1_{\gamma}), \qquad \gamma^x(s^2_{\gamma}) <  \bar \gamma^x (s^2_{\gamma}), \qquad \gamma^v(s^1_{\gamma}, s^2_{\gamma}) \subset A_l.
\end{equation}
Then 
\begin{equation}
    \bs \omega_e(G) = 0.
\end{equation}
\end{lemma}

\begin{thm}\label{thm:besov}
     Let $u$ be a quasi-entropy solution to Burgers equation such that in addition $\mu_1$ given by Proposition \ref{prop:feskin} is a signed measure.
Then for every $T> 0$, $K \subset \mathbb R$ compact and $\delta > 0$ there exists a constant $C \equiv C(K, T, \mu_0)$ such that 
\begin{equation}
\int_{\delta}^T \Vert u(t)\Vert_{B^{1/2, 1}_{\infty}(K)} \dif t < \frac{C}{\min\{\delta,1\}}.
\end{equation}
 \end{thm}

\begin{proof}
Up to a symmetry in the $v$, we can assume that $\mu_1$ is a positive measure.

  Moreover, up to choosing a different pair $(\hat \mu_0, \hat \mu_1)$ smaller than $(\mu_0, \mu_1)$ for $\preceq$, we can assume that $\bs \omega_h$ is a Lagrangian representation of the hypograph of $u$ (Definition~\ref{def:lagrangianrep}) that induces $(\mu_0, \mu_1)$ as in Definition~\ref{defi:goodlr}. Thanks to Proposition~\ref{prop:goodpairs}, we can also assume that $(\mu_0, \mu_1)$ is simultaneously induced (Definition \ref{defi:siminduc}) by $\bs \omega_h$ and $\bs \omega_e$, where $\bs \omega_e$ is a Lagrangian representation of the epigraph of $u$ (Definition~\ref{defi:lagraepi}). This second step is not really necessary but makes the proof easier. A key point is that these operation preserve the sign of $\mu_1$, by definition of the relation $\preceq$. Finally, we denote by $\Gamma_h$, $\Gamma_e$ the set of curves selected by Lemma~\ref{lemma:goodcurves}.
  
\vspace{0.5cm}
\textbf{1.} Fix $\Delta t > 0$ and for $t \geq \Delta t$ consider the set of curves
$$
\Gamma_h^{t, \Delta t} := \Big\{ \gamma_h \in \Gamma_h \quad \big| \quad t_{\gamma_h}^- \leq t-\Delta t, \quad t^+_\gamma \geq t \Big\}
$$
$$
\Gamma_e^{t, \Delta t} := \Big\{ \gamma_e \in \Gamma_e \quad \big| \quad t_{\gamma_e}^- \leq t-\Delta t, \quad t^+_\gamma \geq t \Big\}
$$
Define the measures 
\begin{equation}
    \begin{aligned}
        \chi^t_{a^{\Delta t}} := {e_t}_{\sharp} \;  (\bs \omega_h \llcorner \Gamma_h^{t, \Delta t}) \leq \chi_h(t, \cdot, \cdot )  \msc L^2 \llcorner ( \mathbb R \times (0, 1)) \quad \text{in $\msc M(\mathbb R\times (0, 1))$},& \\
        & \\
         \chi^t_{b^{\Delta t}} := {e_t}_{\sharp} \;  (\bs \omega_e \llcorner \Gamma_e^{t, \Delta t}) \leq \chi_e(t, \cdot, \cdot )  \msc L^2 \llcorner ( \mathbb R \times (0, 1)) \quad \text{in $\msc M(\mathbb R\times (0, 1))$}.
    \end{aligned}
\end{equation}
Finally, we define the functions
\begin{equation}
    \begin{aligned}
       &  a^{\Delta t}(t, x) : = \mr{sup} \; \Big\{ v \in (0, 1) \; \big| \;  (v, x) \in \mr{supp} \; \chi^t_{a^{\Delta t}} \Big\}, \qquad \forall \; x \in \mathbb R\\
       & \\
       &  b^{\Delta t}(t, x) : = \mr{inf} \; \Big\{ v\in (0, 1) \; \big| \;  (v, x) \in \mr{supp} \; \chi^t_{b^{\Delta t}} \Big\}, \qquad \forall \; x \in \mathbb R.
    \end{aligned}
\end{equation}
Notice that the $L^1$ distance between $a^{\Delta t}(t, \cdot)$ and $u(t, \cdot)$ can be estimated in terms of the source $\mu_0$. In fact, by definition of $a^{\Delta t}$, we have
\begin{equation}
\begin{aligned}
    \int_{-M}^M |a^{\Delta t}(t, x) -u(t, x) | \dif x & = \bs \omega_h\Big(\Big\{ \gamma_h \in \Gamma_h \quad \big| \quad t \in I_{\gamma_h}, \quad t^-_{\gamma_h} \in (t-\Delta t, t)\Big\}\Big)\\
    & \\
    & \leq \mu_0^+(S_t^{M, \Delta t})\\
    \end{aligned}
\end{equation}
where 
$$
S_t^{M, \Delta t}:= (t-\Delta t, t) \times (-M- \Delta t, M + \Delta t) \times (0, 1).
$$
This implies, integrating in $(0, T)$ and using Fubini's Theorem:
\begin{equation}\label{eq:uminusa}
    \int_{\Delta t}^T \int_{-M}^M |a^{\Delta t}(t, x) -u(t, x) | \dif x \dif t \leq \Delta t \cdot |\mu_0^+|\big([0,T] \times (-M-\Delta t, M+\Delta t) \times (0, 1) \big).
\end{equation}
Entirely symmetrical statements holds for the distance of $u$ from $b^{\Delta t}$:
\begin{equation}\label{eq:uminusb}
\begin{aligned}
    &  \int_{-M}^M |a^{\Delta t}(t, x) -u(t, x) | \dif x \leq \mu_0^-(S_t^{M, \Delta t}),\\
    & \int_{\Delta t}^T \int_{-M}^M |b^{\Delta t}(t, x) -u(t, x) | \dif x \dif t \leq \Delta t \cdot |\mu_0^-|\big([0,T] \times (-M-\Delta t, M+\Delta t) \times (0, 1) \big).\\
    \end{aligned}
\end{equation}
By triangular inequality, the difference $a^{\Delta t}(t, x)-b^{\Delta t}(t, x)$ lies in $L^1$ as well and 
\begin{equation}\label{eq:abdifference}
\begin{aligned}
& \int_{-M}^M |a^{\Delta t}(t, x)-b^{\Delta t}(t, x)| \dif x  \leq  |\mu_0|(S_t^{M, \Delta t}),\\
    & \int_{\Delta t}^T \int_{-M}^M |a^{\Delta t}(t, x)-b^{\Delta t}(t, x)| \dif x \dif t \leq \Delta t |\mu_0|\big([0,T] \times (-M-\Delta t, M+\Delta t) \times (0, 1) \big).\\
    \end{aligned}
\end{equation}

\vspace{0.5cm}
\textbf{2.}
Now fix $x < y$ and $\varepsilon> 0$ small. Take a curve $\bar \gamma_h \in \Gamma_h^{t, \Delta t}$ such that $|\bar \gamma_h(t) -(y, a^{\Delta t}(t,y))| \leq \varepsilon$. Since $\mu_1$ is positive and $(y, a^{\Delta t}(t,y)) \in \mr{supp} \; \chi_{a^{\Delta t}}^t$, by~\eqref{eq:mu1def} we can assume that $t \mapsto \bar \gamma_h^v(t)$, $t \in I_{\bar \gamma_h}$, is decreasing.
By definition of $b^{\Delta t}$ there exist positive measure set of curves $G \subset \Gamma_e^{t, \Delta t}$ such that  $$|\gamma_e(t)-(x, b^{\Delta t}(t,x))| \leq \varepsilon.$$
By Lemma~\ref{lemma:nocrossing}, it holds
\begin{equation}\label{eq:nocrossingG}
\omega_e\Big( \Big\{ \gamma_e \in G \quad \big| \quad \gamma_e(t-\Delta t) >\bar \gamma_h(t-\Delta t) \quad \text{and} \quad \gamma_e(t) < \bar \gamma_h(t)\Big\}\Big) = 0
\end{equation}
Moreover, since $\mu_1 \geq 0$ and again by \eqref{eq:mu1def}, $I_{\gamma_e} \ni t \mapsto \gamma^v_e(t)$ is increasing for $\omega_e$-a.e. $\gamma_e \in G$. Therefore, thanks to the characteristic equation~\eqref{eq:charaeqla}, it holds
\begin{equation}\label{eq:Oleinikpreineq}
\begin{aligned}
    &\bar \gamma_h^x(t-\Delta t) < y -\Delta t\cdot a^{\Delta t}(t, y) + 2\varepsilon\\
    & \gamma_e^x(t-\Delta t) > x -\Delta t\cdot b^{\Delta t}(t, x) - 2\varepsilon, \quad \text{for $\omega_e$-a.e. $\gamma_e \in G$}.
    \end{aligned}
\end{equation}
Then, combining~\eqref{eq:nocrossingG} with~\eqref{eq:Oleinikpreineq} and letting $\varepsilon \to 0$ we obtain
\begin{equation}
    a^{\Delta t}(t, y) -b^{\Delta t}(t, x) \leq \frac{y-x}{\Delta t} , \qquad \text{for every $x < y$}
\end{equation}
Setting $h = y-x$ and integrating in an interval $[-M, M]$, for some $M > 0$, we obtain
\begin{equation}
    \int_{-M}^M (a^{\Delta t}(t, x+h)-b^{\Delta t}(t, x) )^+ \dif x \leq 2 M h \Delta t^{-1}.
\end{equation}
Moreover, combining with the inequality above the $L^{\infty}$ bounds for $a^{\Delta t}, b^{\Delta t}$,  and the inequality~\eqref{eq:abdifference} of Step 1, we obtain
\begin{equation}
    \begin{aligned}
        -\int_{-M}^M & (a^{\Delta t}(t, x+h)-b^{\Delta t}(t, x))^-  \dif x  = -\int_{-M}^M (a^{\Delta t}(t, x+h)- b^{\Delta t}(t, x))^+ \dif x \\
        &+ \int_{-M}^M (a^{\Delta t}(t, x+h)- b^{\Delta t}(t, x)) \dif x\\
        & \geq -2M h\Delta t^{-1} -2h-|\mu_0|(S_t^{M, \Delta t}) \\
    \end{aligned}
\end{equation}
which in turn yields the bound 
\begin{equation}\label{eq:Oleinik}
\begin{aligned}
    \int_{-M}^M |a^{\Delta t}(t, x+h) &-b^{\Delta t}(t, x)|  \dif x \\
    & \leq C \Big( h\Delta t^{-1}+ |\mu_0|(S_t^{M, \Delta t})\Big),
\end{aligned}
\end{equation}
for all $t > 0$ and $\Delta t < t$, with $C$ depending only on $M$. 

\vspace{0.5cm}
\textbf{Step 3.} Fix $\delta > 0$ and  $h^{1/2} < \delta$. Then, for every $\Delta t < \delta$, we estimate the $L^1$ norm of the difference $\bs u(t,x) -u(t,x+h)$ by
\begin{equation}
    \begin{aligned}
        \int_{\delta}^T \int_{-M}^M |u(t, x) -u(t, x+h)| \dif x  \dif t & \leq  \int_{\delta}^T \int_{-M}^M |u(t, x) -b^{\Delta t}(t, x)| \dif x \dif t+\\
        & + \int_{\delta}^T \int_{-M}^M |b^{\Delta t}(t, x)-a^{\Delta t}(t, x+h)| \dif x \dif t + \\
     & + \int_{\delta}^T\int_{-M}^M |a^{\Delta t}(t, x+h)-u(t, x+h)| \dif x \dif t \\
        & \leq 2 \Delta t\cdot |\mu_0|((0, T) \times (-M-\delta, M +\delta) \times (0, 1)) \\
        & + \int_{\delta}^T \int_{-M}^M |a^{\Delta t}(t, x+h)-b^{\Delta t}(t, x)| \dif x \dif t
    \end{aligned}
\end{equation}
where the inequality in the last line follows by~\eqref{eq:uminusa}, \eqref{eq:uminusb}.
By \eqref{eq:Oleinik}, since for $t \in (\delta, T)$ one has $\Delta t < \delta < t$, we obtain
$$
 \int_{\delta}^T \int _{\mathbb R} |a^{\Delta t}(t, x+h)-b^{\Delta t}(t, x)| \dif x \dif t \leq C(h\Delta t^{-1} + \Delta t)
$$
so that in total, for another constant $C$ depending on $M, T$ and $|\mu_0|$, we obtain
\begin{equation}
    \int_{\delta}^T \int_{-M}^M |u(t, x) -u(t, x+h)| \dif x  \dif t  \leq C(h\Delta t^{-1} + \Delta t), \qquad \forall \; \Delta t < \delta
\end{equation}
Choosing $\Delta t = h^{1/2}$ (which is possible since $h^{1/2} < \delta$ by assumption), we obtain 
\begin{equation}
 \int_{\delta}^T \int_{-M}^M \frac{|u(t, x) -u(t, x+h)|}{h^{1/2}} \dif x  \dif t  <2 C, \qquad \forall h \quad \text{such that} \; h^{1/2} < \delta
\end{equation}
Instead if $h^{1/2} \geq \delta$, we obtain trivially
\begin{equation}
 \int_{\delta}^T \int_{-M}^M \frac{|u(t, x) -u(t, x+h)|}{h^{1/2}} \dif x  \dif t  \leq\frac{2MT}{\delta}.
\end{equation}
This proves that
\begin{equation}
 \sup_{h > 0} \int_{\delta}^T \int_{K} \frac{|u(t, x) -u(t, x+h)|}{h^{1/2}} \dif x  \dif t  \leq \frac{C}{\delta}, \qquad \forall \quad K \subset \mathbb R \; \text{compact}.
\end{equation}
\end{proof}

\end{document}
